\documentclass{easychair}

\usepackage{amsmath,amssymb}
\usepackage{booktabs}
\usepackage{graphicx}
\usepackage{hyperref}
\usepackage{xspace}

\newcommand{\Eprover}{E\xspace}
\newcommand{\Vampire}{Vampire\xspace}
\newcommand{\Mizar}{Mizar\xspace}

\title{Neural Conjecture Generation for Automated Theorem Proving}

\author{
Josef Urban\inst{1}
}

\institute{
Czech Technical University in Prague
}

\authorrunning{Urban}
\titlerunning{Neural Conjecture Generation for ATP}

\begin{document}

\maketitle

\begin{abstract}
We present a neural system for generating intermediate lemmas
(conjectures) that speed up automated theorem provers on large
mathematical libraries.  Given a first-order problem, a graph neural
network encoder builds a heterogeneous representation of the clause
structure, and a Transformer decoder autoregressively generates
candidate cut formulas.  Using the Gentzen cut rule, each candidate is
evaluated by checking that both the problem with the lemma and the
problem with the negated lemma are provable, with proof search savings
measured by the ratio of combined proof effort to the baseline.  We
train and evaluate on the \Mizar Mathematical Library, creating
large-scale datasets with both the \Eprover and \Vampire theorem
provers.  For \Vampire, we construct a dataset of 196\,199 useful
training examples from 2.47 million evaluated lemma--problem pairs
across 34\,182 problems, with proof search reductions of up to $25\times$.
On \Eprover, an ensemble of model variants helps on 125 out of 187
provable test problems (66.8\%), with individual speedups exceeding
$100\times$.  We analyze the complementarity of different model
configurations and provide a budget-optimal strategy for allocating
prover evaluation time across methods.
\end{abstract}

%% ================================================================
\section{Introduction}
\label{sec:intro}

%% TODO: Write introduction - motivation, cut rule, related work, contributions

Automated theorem provers (ATPs) such as \Eprover~\cite{Schulz2002}
and \Vampire~\cite{Kovacs2013} are powerful tools for establishing
logical validity of first-order conjectures.  However, their
performance is often limited by the size of the proof search space,
which can grow exponentially with the number of input clauses.

A classical technique for reducing proof search is the \emph{cut rule}
(Gentzen, 1935): for any clause~$L$, if $P \cup \{L\}$ is
unsatisfiable and $P \cup \{\neg L\}$ is unsatisfiable, then $P$
itself is unsatisfiable.  A well-chosen intermediate lemma~$L$ can
decompose a hard proof into two easier subproofs, dramatically
reducing the total search effort.  The challenge is to find such
lemmas automatically.

We propose a neural approach: a graph neural network (GNN) encodes the
structure of the input problem, and a Transformer decoder generates
candidate lemmas autoregressively.  We train the system on examples
where known lemmas provide proof speedups, and at inference time
generate diverse candidates that are then verified by the ATP.

Our contributions are:
\begin{enumerate}
\item A heterogeneous GNN encoder for CNF problems with five node
  types and role-aware edge convolutions.
\item An autoregressive decoder with arity-constrained generation
  ensuring syntactically valid output.
\item A large-scale \Vampire training dataset with 196K useful
  examples from 34K problems.
\item Comprehensive evaluation showing 66.8\% of provable test
  problems benefit from generated conjectures, with a budget-optimal
  allocation strategy across model variants.
\end{enumerate}


%% ================================================================
\section{Dataset Construction}
\label{sec:dataset}

We create training data by evaluating intermediate lemmas extracted from
successful proofs on the \Mizar Mathematical Library (MML)~\cite{Grabowski2010}.
We describe the two datasets: a smaller one using \Eprover and a larger one
using \Vampire.

\subsection{Problem Corpus}

The \Mizar Mathematical Library contains over 60\,000 theorems
formalized in a rich mathematical language.  We export these as
first-order CNF problems in TPTP format~\cite{Sutcliffe2009}, obtaining
57\,880 problems.

\subsection{Vampire Dataset}
\label{sec:vampire-data}

We use \Vampire with neural clause selection~\cite{Suda2025} to solve the
CNF problems.  Specifically, we run:
\begin{center}\small
\texttt{vampire --decode lrs-1002\_3:2\_av=off:bce=on:...}\\
\texttt{-ncem m3k3kMizChampL17.pt -npcc on -p tptp -i 50000}
\end{center}
with a 50 billion instruction limit.  This solves \textbf{34\,182}
problems (59.1\% of the corpus).

For each solved problem, \Vampire produces proof lemmas---intermediate
formulas that were useful during the proof search.  On average, each
problem yields $\sim$72 such lemmas, for a total of approximately 2.47
million (problem, lemma) pairs.

\paragraph{Cut evaluation.}
For each pair $(P, L)$ where $P$ is a solved problem and $L$ is a
proof lemma, we evaluate the cut rule by running \Vampire twice with
a 25 billion instruction limit:
\begin{itemize}
\item $P \cup \{L\}$: the problem augmented with the lemma as an extra axiom,
\item $P \cup \{\neg L\}$: the problem augmented with the negation of the lemma.
\end{itemize}
If both runs succeed, the \emph{speedup ratio} is
$r = (I_{P \cup L} + I_{P \cup \neg L}) / I_P$,
where $I_X$ denotes the number of instructions burned by \Vampire on
input~$X$.  A ratio $r < 1$ indicates that the lemma provides a
genuine proof search speedup.

\paragraph{Dataset statistics.}
Table~\ref{tab:vampire-stats} summarizes the resulting dataset.
Out of 2\,470\,292 evaluated pairs, 2\,281\,273 have both the lemma
and its negation proved within the 25B instruction limit.  Of these,
\textbf{196\,199} provide a speedup ($r < 1$), covering 3\,591
distinct problems.

\begin{table}[t]
\centering
\caption{Vampire conjecture dataset statistics.}
\label{tab:vampire-stats}
\begin{tabular}{lr}
\toprule
CNF problems attempted & 57\,880 \\
Solved by Vampire (50B instr.) & 34\,182 \\
(Problem, lemma) pairs evaluated & 2\,470\,292 \\
Both directions proved (25B instr.) & 2\,281\,273 \\
Useful speedups ($r < 1$) & 196\,199 \\
Problems with $\geq 1$ useful lemma & 3\,591 \\
\bottomrule
\end{tabular}
\end{table}

Table~\ref{tab:ratio-dist} shows the distribution of speedup ratios.
A substantial fraction of useful examples provide large speedups:
3\,190 examples (from 164 problems) achieve a ratio $\leq 0.1$,
corresponding to a $10\times$ or greater reduction in proof search
effort.  The median useful ratio is approximately 0.7, and the dataset
contains examples across the full range of speedup magnitudes.

\begin{table}[t]
\centering
\caption{Distribution of useful examples by speedup ratio threshold.}
\label{tab:ratio-dist}
\begin{tabular}{rrrrr}
\toprule
Ratio $\leq$ & Examples & Problems & \% of useful & \% of problems \\
\midrule
0.05 &        83 &    10 &  0.0\% &  0.3\% \\
0.10 &    3\,190 &   164 &  1.6\% &  4.6\% \\
0.20 &   15\,165 &   553 &  7.7\% & 15.4\% \\
0.30 &   30\,933 &   955 & 15.8\% & 26.6\% \\
0.50 &   66\,461 & 1\,645 & 33.9\% & 45.8\% \\
0.70 &  110\,395 & 2\,366 & 56.3\% & 65.9\% \\
0.90 &  163\,531 & 3\,156 & 83.3\% & 87.9\% \\
0.95 &  179\,059 & 3\,371 & 91.3\% & 93.9\% \\
1.00 &  196\,199 & 3\,591 & 100\% & 100\% \\
\bottomrule
\end{tabular}
\end{table}

\paragraph{Data imbalance.}
The useful examples are unevenly distributed across problems.  Some
problems have over 700 useful lemmas (e.g., \texttt{t3\_gate\_2} with
718), while 143 problems have exactly one.  The mean is 54.6 useful
lemmas per problem and the median is 35.  Table~\ref{tab:imbalance}
shows the cumulative distribution.  This imbalance must be addressed
during training to prevent the model from overfitting to
heavily-represented problems.  We use per-problem weighting: each
problem contributes equally to the loss regardless of how many
training examples it provides.

\begin{table}[t]
\centering
\caption{Data imbalance: distribution of useful examples per problem.}
\label{tab:imbalance}
\begin{tabular}{rrrr}
\toprule
$\leq k$ useful & Problems & Cumul. \% \\
\midrule
1   &   143 &  4.0\% \\
5   &   470 & 13.1\% \\
10  &   758 & 21.1\% \\
20  & 1\,231 & 34.3\% \\
50  & 2\,243 & 62.5\% \\
100 & 3\,018 & 84.0\% \\
200 & 3\,480 & 96.9\% \\
500 & 3\,586 & 99.9\% \\
\bottomrule
\end{tabular}
\end{table}


\subsection{E Prover Dataset}
\label{sec:eprover-data}

For initial development and comparison, we also create a smaller
dataset using the \Eprover theorem prover.  From a subset of 3\,161
CNF problems, we obtain 122\,356 training examples with a
train/validation/test split of 2\,239/279/279 problems.  The
evaluation uses processed clauses as the cost metric (rather than
instructions), with a 10-second wall-clock timeout.


%% ================================================================
\section{Architecture}
\label{sec:arch}

%% TODO: GNN encoder, Transformer decoder, arity constraints, pointer mechanism

\subsection{Heterogeneous GNN Encoder}

\subsection{Autoregressive Decoder}

\subsection{Arity-Constrained Generation}

\subsection{Model Variants}


%% ================================================================
\section{Training}
\label{sec:training}

%% TODO: Loss function, per-problem weighting, data augmentation


%% ================================================================
\section{Evaluation}
\label{sec:eval}

%% TODO: E prover results, Vampire results, complementarity, budget analysis


%% ================================================================
\section{Related Work}
\label{sec:related}

%% TODO: Cut introduction in ATPs, neural theorem proving, lemma generation


%% ================================================================
\section{Conclusion}
\label{sec:conclusion}

%% TODO


%% ================================================================
\bibliographystyle{plain}
\begin{thebibliography}{10}

\bibitem{Schulz2002}
S.~Schulz.
\newblock E -- a brainiac theorem prover.
\newblock {\em AI Communications}, 15(2-3):111--126, 2002.

\bibitem{Kovacs2013}
L.~Kov{\'a}cs and A.~Voronkov.
\newblock First-order theorem proving and {V}ampire.
\newblock In {\em CAV 2013}, LNCS 8044, pp.~1--35, 2013.

\bibitem{Suda2025}
M.~Suda.
\newblock Efficient neural clause-selection for {V}ampire.
\newblock {\em arXiv:2503.07792}, 2025.

\bibitem{Grabowski2010}
A.~Grabowski, A.~Korni{\l}owicz, and A.~Naumowicz.
\newblock Mizar in a nutshell.
\newblock {\em J. Formalized Reasoning}, 3(2):153--245, 2010.

\bibitem{Sutcliffe2009}
G.~Sutcliffe.
\newblock The {TPTP} problem library and associated infrastructure.
\newblock {\em J. Automated Reasoning}, 43(4):337--362, 2009.

\end{thebibliography}

\end{document}
